
\section{Computation of integrals over four-dimensional regions}

\SourcePage{651}{656}

We shall present here certain rules and formulae useful for computing the integrals that arise in the theory of radiative corrections. The typical formula for an integral corresponding to a Feynman diagram is
\begin{equation}
\int \frac{f(k)\,d^4k}{a_1 a_2 \ldots a_n},
\label{eq:131.1}
\end{equation}
where $a_1, a_2, \ldots$ are polynomials of second degree in the 4-vector $k$, $f(k)$ is a polynomial of some degree $n'$, and the integration is taken over the entire four-dimensional $k$-space.

A convenient method of computing such integrals (due to \emph{Feynman}, 1949) is based on a preliminary transformation (parametrisation) of the integrand by introducing additional integrations over auxiliary variables $\xi_1, \xi_2, \ldots$ according to the formula
\begin{equation}
\frac{1}{a_1 a_2 \ldots a_n} = (n-1)! \int_0^1 d\xi_1 \int_0^1 \ldots \int_0^1 d\xi_n \,\frac{\delta(\xi_1 + \xi_2 + \ldots + \xi_n - 1)}{(a_1\xi_1 + a_2\xi_2 + \ldots + a_n\xi_n)^n}.
\label{eq:131.2}
\end{equation}
As a result of such a transformation, instead of $n$ distinct quadratic polynomials in the denominator there arises a degree $n$ of just a single polynomial of second degree.

Eliminating the $\delta$-function by integrating over $d\xi_n$ and introducing new variables according to
\[
\xi_1 = x_{n-1},\quad \xi_2 = x_{n-2} - x_{n-1},\quad \ldots, \quad \xi_{n-1} = x_1 - x_2,
\]
\[
\xi_1 + \xi_2 + \ldots + \xi_{n-1} = x_1,
\]
we obtain formula (\ref{eq:131.2}) in the equivalent form:
\begin{equation}
\frac{1}{a_1 a_2 \ldots a_n} = (n-1)! \int_0^1 dx_1 \int_0^{x_1} dx_2 \ldots \int_0^{x_{n-2}} dx_{n-1}\,\frac{1}{[a_1 x_{n-1} + a_2(x_{n-2} - x_{n-1}) + \ldots + a_n(1 - x_1)]^n}.
\label{eq:131.3}
\end{equation}

\SourcePage{652}{657}

For $n = 2$ this formula takes the form
\begin{equation}
\frac{1}{a_1 a_2} = \int_0^1 \frac{dx}{[a_1 x + a_2(1-x)]^2}
\label{eq:131.4}
\end{equation}
and is verified by direct computation. For arbitrary $n$ the formula can be proved by induction from $n - 1$ to $n$. Indeed, carrying out in (\ref{eq:131.3}) the integration over $dx_{n-1}$, one obtains on the right-hand side the difference of two $(n-2)$-fold integrals of the same form. Assuming the formula to be valid for them, we get $\dfrac{1}{a_1 - a_2}\biggl[\dfrac{1}{a_2 a_3 \ldots a_n} - \dfrac{1}{a_1 a_3 \ldots a_n}\biggr]$, which coincides with the expression on the left-hand side of (\ref{eq:131.3}).

By differentiating (\ref{eq:131.3}) with respect to $a_1, a_2, \ldots$ one can obtain analogous formulae serving for the parametrisation of integrals containing in the denominator any of the polynomials raised to powers higher than the first.

Regularisation of divergent integrals is effected by subtracting from them integrals of analogous form. For the computation of such a difference it may prove expedient to carry out a preliminary transformation of the difference of the integrands (each of which has already been transformed by means of (\ref{eq:131.2})) with the aid of the formula
\begin{equation}
\frac{1}{a^n} - \frac{1}{b^n} = -\int_0^1 \frac{n(a-b)\,dz}{[(a-b)z + b]^{n+1}}.
\label{eq:131.5}
\end{equation}

After the transformation, according to (\ref{eq:131.3}), the four-dimensional integration in (\ref{eq:131.1}) reduces to the form
\begin{equation}
\int \frac{f(k)\,d^4k}{[(k - l)^2 - \alpha^2]^n},
\label{eq:131.6}
\end{equation}
where $l$ is a 4-vector and $\alpha^2$ a scalar, both depending on the parameters $x_1, \ldots, x_{n-1}$; the scalar $\alpha^2$ will be taken to be positive.

If the integral (\ref{eq:131.6}) converges, then one can carry out the substitution $k - l \to k$ (shift of the origin of coordinates), after which it takes the form
\begin{equation}
\int \frac{f(k)\,d^4k}{(k^2 - \alpha^2)^n}
\label{eq:131.7}
\end{equation}
(with a different function $f(k)$), so that the denominator contains only the square $k^2$. As regards the numerator, it suffices to

\SourcePage{653}{658}

\noindent consider scalar functions $f = F(k^2)$. Indeed, for integrals of other types with respect to the numerator we have
\begin{equation}
\int \frac{k^\mu\,F(k^2)\,d^4k}{(k^2 - \alpha^2)^n} = 0,
\label{eq:131.8}
\end{equation}
\begin{equation}
\int \frac{k^\mu k^\nu\,F(k^2)\,d^4k}{(k^2 - \alpha^2)^n} = \frac{1}{4}\,g^{\mu\nu}\int \frac{k^2\,F(k^2)\,d^4k}{(k^2 - \alpha^2)^n},
\label{eq:131.9}
\end{equation}
\begin{multline}
\int \frac{k^\mu k^\nu k^\rho k^\sigma\,F(k^2)\,d^4k}{(k^2 - \alpha^2)^n} \\
= \frac{1}{24}\bigl(g^{\mu\nu}g^{\rho\sigma} + g^{\mu\rho}g^{\nu\sigma} + g^{\mu\sigma}g^{\nu\rho}\bigr)\int \frac{(k^2)^2\,F(k^2)\,d^4k}{(k^2 - \alpha^2)^n},
\label{eq:131.10}
\end{multline}
and so on, which is evident from symmetry considerations (upon integration over all directions of $k$).

In the original integral (\ref{eq:131.1}), each of the factors $a_1, a_2, \ldots$ in the denominator has (as a function of $k_0$) two zeros, which are dealt with upon integration over $dk_0$ according to the usual rule (see \S\,75). After the transformation to the form (\ref{eq:131.7}), instead of $2n$ simple poles the integrand has only two poles of order $n$, which are dealt with by the same rule (contour $C$ in Fig.~25). Rotating the contour of integration, as indicated by the arrows in the figure, one can bring it into coincidence with the imaginary axis in the $k_0$-plane ($C'$ in Fig.~25). In other words, the variable $k_0$ is replaced by $k_0 =$

% Figure 25: Contour rotation in the complex k_0 plane.
% The contour C runs along the real axis, bypassing poles according to the Feynman rule.
% The contour C' runs along the imaginary axis after rotation.

\noindent $= ik'_0$ with a real variable $k'_0$. Changing also the notation $\mathbf{k}$ to $\mathbf{k}'$, we have
\begin{equation}
k^2 = k_0^2 - \mathbf{k}^2 \to -(k_0'^2 + \mathbf{k}'^2) = -k'^2,
\label{eq:131.11}
\end{equation}
where $k'$ is a 4-vector in Euclidean metric. In this case
\[
d^4k \to i\,d^4k' = i\,k'^2\,dk'^2\,d\Omega,
\]
where $d\Omega$ is an element of four-dimensional solid angle. Integration over $d\Omega$ gives $2\pi^2$ (see II, \S\,111), after which
\begin{equation}
d^4k \to i\pi^2\,k'^2\,d(k'^2).
\label{eq:131.12}
\end{equation}
Denoting $k'^2 = z$, we obtain finally

\SourcePage{654}{659}

\begin{equation}
\int \frac{F(k^2)\,d^4k}{(k^2 - \alpha^2)^n} = (-1)^n\,i\pi^2 \int_0^\infty \frac{F(-z)\,z\,dz}{(z + \alpha^2)^n}.
\label{eq:131.13}
\end{equation}

In particular,
\begin{equation}
\int \frac{d^4k}{(k^2 - \alpha^2)^n} = \frac{(-1)^n\,i\pi^2}{\alpha^{2(n-2)}(n-1)(n-2)}.
\label{eq:131.14}
\end{equation}

The logarithmically divergent part in integrals of the form (\ref{eq:131.7}) can be separated out as
\begin{equation}
\int \frac{d^4k}{[(k-l)^2 - \alpha^2]^2}.
\label{eq:131.15}
\end{equation}
It is easy to see that in such an integral the substitution $k \to k + l$ is permissible. Indeed, the difference of the original and the transformed integrals
\[
\int \biggl\{\frac{1}{[(k-l)^2 - \alpha^2]^2} - \frac{1}{(k^2 - \alpha^2)^2}\biggr\}\,d^4k
\]
represents a convergent integral, and hence the replacement $k \to k + l$ within it is permissible in any case. Carrying it out and then replacing $k \to -k$, we obtain the same quantity with the opposite sign, whence its vanishing follows.

A linearly divergent integral must have the form
\begin{equation}
\int \frac{k^\mu\,d^4k}{[(k-l)^2 - \alpha^2]^2},
\label{eq:131.16}
\end{equation}
but in fact such an integral diverges only logarithmically: the integrand asymptotically (as $k \to \infty$) equals $k^\mu/(k^2)^2$ and vanishes upon averaging over directions. The shift of the origin of coordinates, however, does not leave the integral (\ref{eq:131.16}) unchanged, but adds to it an additive constant. We shall demonstrate this for the case of an infinitesimally small shift $k \to k + \delta l$, computing the difference
\begin{equation}
\Delta^\mu = \int \biggl\{\frac{k^\mu}{[(k - \delta l)^2 - \alpha^2]^2} - \frac{k^\mu + \delta l^\mu}{(k^2 - \alpha^2)^2}\biggr\}\,d^4k.
\label{eq:131.17}
\end{equation}
To the accuracy of terms of first order in $\delta l$,
\[
\Delta^\mu = \int \biggl\{\frac{4k^\mu(k\delta l)}{(k^2 - \alpha^2)^3} - \frac{\delta l^\mu}{(k^2 - \alpha^2)^2}\biggr\}\,d^4k.
\]
In the first term the averaging over directions replaces the numerator by $k^2\,\delta l^\mu$ (cf.\ (\ref{eq:131.9})), after which we find\footnote{A more laborious calculation leads to the same result also for finite $l$.}
\begin{equation}
\Delta^\mu = \alpha^2\,\delta l^\mu \int \frac{d^4k}{(k^2 - \alpha^2)^3} = -\frac{i\pi^2}{2}\,\delta l^\mu.
\label{eq:131.18}
\end{equation}

\SourcePage{655}{660}

In the final expressions for radiative corrections there frequently appears a transcendental function defined by the integral
\begin{equation}
F(\xi) = \int_0^\xi \frac{\ln(1+x)}{x}\,dx
\label{eq:131.19}
\end{equation}
(it is sometimes called the \emph{Spence function}). We shall note here for reference some of its properties:
\begin{equation}
F(\xi) + F\!\left(\frac{1}{\xi}\right) = \frac{\pi^2}{6} + \frac{1}{2}\ln^2\xi,
\label{eq:131.20}
\end{equation}
\begin{equation}
F(-\xi) + F(-1+\xi) = -\frac{\pi^2}{6} + \ln\xi\,\ln(1 - \xi),
\label{eq:131.21}
\end{equation}
\begin{equation}
F(1) = \frac{\pi^2}{12}, \qquad F(-1) = -\frac{\pi^2}{6}.
\label{eq:131.22}
\end{equation}
Expansion for small $\xi$:
\begin{equation}
F(\xi) = \xi - \frac{\xi^2}{4} + \frac{\xi^3}{9} - \frac{\xi^4}{16} + \ldots
\label{eq:131.23}
\end{equation}
